\documentclass[11pt,a4paper]{article}

\usepackage[utf8]{inputenc}
\usepackage[T2A]{fontenc}
\usepackage[russian,english]{babel}
\usepackage{geometry}
\usepackage{array}
\usepackage{hyperref}
\usepackage{xcolor}

\geometry{left=17mm,right=17mm,top=15mm,bottom=17mm}
\hypersetup{
    colorlinks=true,
    linkcolor=blue!60!black,
    urlcolor=blue!60!black,
    citecolor=blue!60!black
}
\setlength{\parindent}{0pt}
\setlength{\parskip}{4pt}

\begin{document}
\selectlanguage{russian}

\begin{center}
{\Large\bfseries Complexity-aware KAN для нелинейной эквализации 16QAM-сигнала}\\
\vspace{2mm}
Направление AIRI: \textbf{EFFICIENTDL} --- эффективный ИИ: оптимизация вычислений, ускорение инференса и системный дизайн современных моделей
\end{center}

\textbf{Выбранная статья.}
В качестве базовой работы выбрана статья \textit{KAN: Kolmogorov--Arnold Networks} Z.~Liu, Y.~Wang, S.~Vaidya, F.~Ruehle, J.~Halverson, M.~Soljačić, T.~Y.~Hou, M.~Tegmark, опубликованная на ICLR 2025 Oral \cite{kan}. ICLR имеет ранг CORE A* в CORE2023 \cite{core}. Работа предлагает Kolmogorov--Arnold Networks (KAN) как альтернативу MLP: вместо фиксированных нелинейностей в узлах и обучаемых линейных весов на ребрах KAN использует обучаемые одномерные функции на ребрах, обычно параметризованные B-сплайнами. Это делает архитектуру привлекательной для задач, где требуется компактно аппроксимировать гладкие нелинейные зависимости.

\textbf{Анализ метода и ключевые задачи.}
Основная идея KAN состоит в замене слоя вида \(x \mapsto \sigma(Wx)\) на композицию матриц обучаемых одномерных функций. Авторы мотивируют это теоремой Колмогорова--Арнольда и показывают, что KAN может давать лучшую интерпретируемость, масштабирование и качество на ряде научных задач. Ключевые задачи метода: (1) выбрать параметризацию функций на ребрах; (2) обеспечить устойчивое обучение сетки и коэффициентов сплайна; (3) сохранить выигрыш качества без чрезмерного роста памяти и времени инференса.

Сильная сторона статьи --- ясная архитектурная альтернатива MLP: KAN лучше соответствует задачам аппроксимации гладких нелинейных функций и потенциально допускает визуализацию/разреживание ребер. Слабая сторона --- вычислительная цена и чувствительность к гиперпараметрам: размер сетки, порядок сплайна, ширина и число слоев могут менять как BER/ошибку, так и скорость на порядки. Кроме того, в статье недостаточно изучен режим, где метрика качества не является MSE, а задается прикладной дискретной метрикой, например BER в задачах связи.

\textbf{Предлагаемое развитие.}
Я предлагаю развить KAN в направлении \textit{complexity-aware nonlinear equalization}: применить KAN к посимвольной эквализации 16QAM-сигнала и изучить компромисс между BER, числом параметров и скоростью. Вход модели --- окно принятых IQ-символов длины \(65\), выход --- либо исправленные координаты \((I,Q)\), либо один из 16 классов созвездия. Такая постановка проверяет, действительно ли обучаемые одномерные нелинейности полезны в физически мотивированной задаче, где канал вносит гладкие, но сложные искажения. Практическое улучшение статьи: добавить BER-aware selection of KAN complexity, pruning/FastKAN-варианты и гибридную функцию потерь, совмещающую регрессию \((I,Q)\) с символной ошибкой.

\textbf{Минимальные эксперименты.}
Были проведены пилотные эксперименты на датасете 16QAM: train-файлы 1--56, validation 57--62, test 63--64; тестовые файлы не используются при обучении, выборе чекпоинта и нормализации. Базовый BER без нейросетевой эквализации: \(9.71\cdot 10^{-3}\). Сравнивались MLP, KAN-регрессор и KAN-классификатор. BER считается после отображения предсказания в ближайшую точку 16QAM-созвездия и сравнения Gray-битов.

\begin{center}
\small
\begin{tabular}{|>{\raggedright\arraybackslash}p{38mm}|r|r|r|r|}
\hline
\textbf{Модель / sweep} & \textbf{Параметры} & \textbf{BER test} & \textbf{Улучшение} & \textbf{Скорость, sym/s} \\
\hline
MLP-регрессор & 27k & \(2.27\cdot10^{-3}\) & 76.7\% & \(3.8\cdot10^{6}\) \\
\hline
KAN-регрессор, 250 эпох & 832k & \(1.93\cdot10^{-3}\) & 80.1\% & \(1.63\cdot10^{5}\) \\
\hline
KAN-классификатор, 250 эпох & 877k & \(\mathbf{8.90\cdot10^{-4}}\) & \(\mathbf{90.8\%}\) & \(1.61\cdot10^{5}\) \\
\hline
KAN-cls grid \(4 \to 16\), 60 эпох & 316k--737k & \(2.17\cdot10^{-3}\to\mathbf{8.49\cdot10^{-4}}\) & 77.7--91.3\% & \(4.20\cdot10^{5}\to1.88\cdot10^{5}\) \\
\hline
KAN-cls spline order \(1,2,3,4\), 60 эпох & 386k--491k & \(1.51,1.41,\mathbf{1.39},1.48\)\(\cdot10^{-3}\) & 84.5--85.7\% & \(1.02\cdot10^{6}\to2.03\cdot10^{5}\) \\
\hline
KAN-cls hidden \(64\to192\), grid 8, 60 эпох & 175k--844k & \(2.06\cdot10^{-3}\to1.15\cdot10^{-3}\) & 78.9--88.2\% & \(4.86\cdot10^{5}\to2.15\cdot10^{5}\) \\
\hline
\end{tabular}
\end{center}

Результаты показывают два вывода. Во-первых, KAN-классификатор существенно снижает BER относительно MLP и KAN-регрессора: лучший запуск дает \(8.90\cdot10^{-4}\), то есть примерно в 10.9 раза ниже исходного BER. Во-вторых, выигрыш достигается высокой ценой: MLP примерно в 20--24 раза быстрее и на порядок компактнее. Поэтому наиболее перспективная линия развития --- не просто ``заменить MLP на KAN'', а построить компактный KAN-эквалайзер: подобрать grid около 16, проверить hidden-dim sweep при фиксированном grid=16, затем применить pruning или RBF/FastKAN-базис и сравнить точку на кривой BER--complexity.

\textbf{Итог.}
Статья KAN является сильной Core A* основой для направления EFFICIENTDL, а предложенное развитие переводит ее из общей архитектурной идеи в прикладную задачу эффективной связи. Минимальные эксперименты подтверждают потенциал KAN по BER, но также выявляют главный исследовательский вопрос: как сохранить выигрыш KAN-классификатора при сложности, близкой к MLP.

\begin{thebibliography}{9}
\bibitem{kan}
Z. Liu, Y. Wang, S. Vaidya, F. Ruehle, J. Halverson, M. Soljačić, T. Y. Hou, M. Tegmark.
\textit{KAN: Kolmogorov--Arnold Networks}. ICLR 2025 Oral.
\url{https://openreview.net/forum?id=Ozo7qJ5vZi}

\bibitem{core}
CORE Conference Ranking Portal, CORE2023: International Conference on Learning Representations (ICLR), rank A*.
\url{https://portal.core.edu.au/conf-ranks/?by=all&page=1&sort=arank&source=CORE2023}
\end{thebibliography}

\end{document}
